\scanpage{57}
\mspage{11}
$\underline C$ with a ACU constraint $(a,c,(\underline{1},g,d))$, a $\otimes$-category $\underline C'$ with a ACU constraint $(a',c',(\underline{1}',g',d'))$, the base category of which is a groupoïd, and a $\otimes$-ACU functor $(F,\breve F):\underline C'\longrightarrow\underline C$. We try to find a $\otimes$-ACU category $\underline P$ and a $\otimes$-ACU functor $(\unsure{\mathfrak D|\mathscr D|Q},\breve{\mathfrak D}):\underline C\longrightarrow\underline P$ having the following properties:

$1^\circ$ $\mathfrak DFX'$ is invertible in $\underline P$ for all $X'\in\operatorname{Ob}\underline C'$.

$2^\circ$ For all $\otimes$-ACU functor $(\unsure{\mathfrak E|\mathscr E|Y},\breve{\mathfrak E})$ from $\underline C$ to a $\otimes$-ACU category $\underline Q$
% slanted
\margin{\unsure{give} the \unsure{equivalence|equivalency|équivalence} of categories \tikz[baseline]{\draw (0,.55) -- (.025,-.28) .. controls (.025,-.31) and (-.015,-.31) .. (-.025,-.28);}}
such that $\mathfrak EFX'$ is invertible in $\underline Q$ for all $X'\in\operatorname{Ob}\underline C'$, there exists a $\otimes$-ACU functor $(E',\breve E')$, unique up to $\otimes$-isomorphism, from $\underline P$ to $\underline Q$ such that $(\mathfrak E,\breve{\mathfrak E})\simeq(E',\breve E')\circ(\mathfrak D,\breve{\mathfrak D})$.

This problem is reduced by the first by putting $\underline A'=\underline C'$, $\underline A=\underline C\times\underline C'$, $TX'=(FX',X')$ and by remarking that if $\underline C,\underline C',\underline Q$ are $\otimes$-ACU categories, $\underline{\operatorname{Hom}}^{\otimes,\mathrm{ACU}}(\underline C,\underline Q)$ the category of all $\otimes$-ACU functor from $\underline C$ to $\underline Q$, then there is a canonical equivalence of categories
\[
\resizebox{\linewidth}{!}{$\displaystyle
\underline{\operatorname{Hom}}^{\otimes,\mathrm{ACU}}(\underline C\deleted{\unsure{\otimes|\times}}\unsure{\times|\otimes}\underline C',\underline Q)\longrightarrow
\underline{\operatorname{Hom}}^{\otimes,\mathrm{ACU}}(\underline C,\underline Q)\times
\underline{\operatorname{Hom}}^{\otimes,\mathrm{ACU}}(\underline C',\underline Q)$}
\]
The $\otimes$-ACU category $\underline P$ thus defined is called the \uline{$\otimes$-category of fractions} of the category $\underline C$ with respect to $(\underline C',(F,\breve F))$. The $\otimes$-category of fractions of $\underline C^{\mathrm{is}}$ with respect to $(\underline C^{\mathrm{is}},(\mathrm{id}_{\underline C^{\mathrm{is}}},\mathrm{id}))$ is a Pic-category which is called the \uline{Pic-envelope} of the category $\underline C$, and denoted by $\operatorname{Pic}(\underline C)$.

For an application of the Pic-envelope, we take $\underline C=\unsure{\mathcal P|\mathscr P}(R)$, category of all finitely generated $R$-modules ($R$ a ring) and $\underline P=\operatorname{Pic}(\mathcal P(R))$, then one obtains
\[
\begin{aligned}
\Pi_0(\underline P)&\simeq K^0(R)\\
\Pi_1(\underline P)&\simeq K^1(R)
\end{aligned}
\]
% slanted
\margin{this Pic-envelope is \uline{not} a \uline{strict} Pic-category}
where $K^0(R)$ is the Grothendieck group and $K^1(R)$ the Whitehead group [1]

The use of the $\otimes$-category of fractions of a $\otimes$-ACU category gives us the following result:
% archive sheet number: 56
